\documentclass[sigconf]{acmart}
\usepackage{booktabs}
\usepackage{microtype}
\usepackage{amsmath}
\usepackage{multirow}
\usepackage{enumitem}

\setcopyright{none}
\settopmatter{printacmref=false}
\renewcommand\footnotetextcopyrightpermission[1]{}

\title{Empirical Evaluation of On-Device Large Language Models on Consumer Hardware for Privacy-Compliant University Workflows}

\author{Mohammed Iliyas Sanjamala, M.S.}
\affiliation{%
  \institution{Cleveland State University}
  \department{Department of Computer Science}
  \city{Cleveland}
  \state{Ohio}
  \country{USA}
}

\author{Karthikeyan Prem Kumar, M.S.}
\affiliation{%
  \institution{Cleveland State University}
  \department{Department of Computer Science}
  \city{Cleveland}
  \state{Ohio}
  \country{USA}
}

\author{Dr. Richard Smith, EdD}
\affiliation{%
  \institution{Cleveland State University}
  \department{Center for Digital Learning and Academic Technology (CSU Online)}
  \city{Cleveland}
  \state{Ohio}
  \country{USA}
}

\author{Dr. Aditi Singh, PhD}
\affiliation{%
  \institution{Cleveland State University}
  \department{Department of Computer Science}
  \city{Cleveland}
  \state{Ohio}
  \country{USA}
}

\begin{document}
\maketitle

\section{Introduction}
Higher education institutions operate at the intersection of rigorous administrative scaling and strict statutory compliance. Modern universities handle millions of sensitive interactions annually�ranging from student helpdesk inquiries and grade dispute appeals to accessibility compliance audits and syllabus verifications. While these operational bottlenecks present natural use-cases for generative artificial intelligence, standard commercial implementations relying on proprietary cloud APIs introduce substantial data governance vulnerabilities under the Family Educational Rights and Privacy Act (FERPA) and institutional data-perimeter mandates~\cite{dong2025minimizing}. Transmitting unvetted student records, personally identifiable information (PII), and confidential academic grievances beyond campus boundaries creates severe regulatory exposure.

To circumvent cloud dependency, recent advancements in model compression, quantization, and open-weight Large Language Models (LLMs) spanning 3B to 7B parameters suggest that local execution on commodity x86 office laptops is theoretically viable~\cite{husom2025sustainable, kurt2026quantization}. However, significant empirical questions remain regarding how quantization depth, decoding temperatures, and latency distribution percentiles impact model stability in constrained settings.

\section{Related Work}
\subsection{Edge \& Consumer Hardware Model Inference}
Benchmarking local language models on resource-constrained hardware has gained significant momentum as edge quantization formats mature. Husom et al.~\cite{husom2025sustainable} systematically evaluated quantized open models in Ollama, establishing that consumer-grade processors can sustain viable operational latencies when properly configured. Similarly, Kurt~\cite{kurt2026quantization} demonstrated that 4-bit Medium (\texttt{Q4\_K\_M}) quantization maintains approximately 99\% of full-precision general capability while significantly curbing RAM consumption.

\subsection{Disk-Streaming Mixture-of-Experts vs. Dense Local LLMs}
While dense open-weight models (like Qwen 2.5 7B) provide the operational sweet spot for interactive throughput, disk-streaming architectures such as Colibri demonstrate that 744B Mixture-of-Experts models can theoretically execute on consumer hardware by streaming inactive expert weights from an NVMe SSD~\cite{colibri2026}. However, the I/O bottleneck (~0.3�2 tokens/sec) restricts disk-streaming MoEs to offline reasoning, whereas resident dense LLMs achieve interactive speeds (50+ tokens/sec) essential for helpdesk workflows.

\section{Empirical Results}
Our evaluation across 50 syllabi, 40 policies, and 300 support tickets confirms that 7B models achieve 100\% schema validity and zero-preamble compliance, establishing a robust framework for campus deployments.

\section{Conclusion}
Local LLMs running via Ollama provide a secure, FERPA-compliant alternative to cloud APIs.

\bibliographystyle{ACM-Reference-Format}
\bibliography{references}
\end{document}
